\section{A three-letter arc with a separated middle up letter}

% BEGIN F83S1
\paragraph{F83S1. A short visible coefficient despite the initial-index factor.}
Take a length-three mixed arc in U83M1, in forward order down, up,
down. Put $y=s_{c-2}+S$, $a=s_{c-1}$, $z=y+\eta$.
Consider the original separated case $pz\le a$ including equality.
Here $a\in[B,B+\eta]$, $20\eta,20S,20w\le B$, $h=B+\eta+S$,
and $40U\le B/4$. Write $I=\mathrm{AI}(F,c-1)$, as in
S83L3. R82W1 gives
$I\le(h+\eta)/(p-1)+1$, using the common upper bound $h+\eta$
for the signed interval widths. Set $Q=v+p$, $P=P_1$, and
$t=\lfloor3U\rfloor+1$. There exists $j<a-\eta$ such that the
coefficient of $Q^j$ in
\[
                          (Q-p)^I P^z
\]
has normalized $p$-valuation less than $t$.

Indeed $I+t+\eta<a$. To see this, use $p-1\ge10$,
$h+\eta\le23B/20$, $t\le3B/160+1$, and $\eta\le B/20$;
$B\ge160$ follows from $U\ge1$. If all the indicated coefficients
were divisible by $p^t$, then
$(Q-p)^IP^z\in(p^t,Q^{a-\eta})$ in $O[[Q]]$.
Expand $Q^{I+t-1}=((Q-p)+p)^{I+t-1}$. Every term is divisible
by $(Q-p)^I$ or by $p^t$. Hence
$Q^{I+t-1}P^z\in(p^t,Q^{a-\eta})$ and every coefficient of
$P^z$ below degree $D=a-\eta-I-t+1>0$ is divisible by $p^t$.
Use $P=p\varepsilon+Q^p V$ with $V,\varepsilon$ units as in
S83L4. For any $k\le z$ with $pk<D$ it follows successively that
\[
                       t\le z-k+v_p\binom zk .
\]
Indeed subtract the previous integral multiples of $p^t$ in
the binomial expansion and read the term in degree $pk$, whose
unit factor has nonzero residue. If $pz<D$, take $k=z$ to get
$t\le0$. Otherwise take $k=\lfloor(D-1)/p\rfloor$ and use the
factorial bound of S83L4. Since $z\le a/p$, this gives
\[
  t\le z-D/p+1+\log_p z
        \le (\eta+I+t)/p+1+\log_p(2h).
\]
But $(I+\eta)/p\le(11/10)(h/p^2+\eta/p)+1/p$ using $p\ge11$.
Adding $1+\log_p(2h)$ yields at most $(11/10)U$ since
$1+1/p\le11/10$. In contrast $t(1-1/p)>(30/11)U$.
This contradiction proves the coefficient assertion. The argument
uses powers of the rational prime in the original, possibly
ramified, coefficient DVR.
% END F83S1

% BEGIN F83S2
\paragraph{F83S2. A filtered direction for the full down, up, down scalar test.}
Use the first two forward steps, from the preceding cut $c-3$ to the
middle letter $c-1$, keeping the target direction $E_2$ of S83L1.
It belongs to $\mathrm{Fil}^a$ at this endpoint, after transporting
back by one crystalline Frobenius. On these steps take the test
$P^yE_2/Q^a$, with $P,Q$ as in F83S1. Its numerator denominator
bound is $\Lambda=Q^hP^G$, $G=w+S$. The test passes the true
target filtration at both levels since the first down letter has
minimum jump $-y$. U83M6's direct comparison gives an element
of $L_{2,c-1}$ representing $p^\sigma$ times this test.
Here only the initial part of the three-letter arc is used.
Indeed the same predecessor basis is carried through two Frobenius
steps and the modulus is $\varphi(v^MP_0^{2G})$, of orders
$qM$ at $Q$ and $2G$ at $P$. Thus all frame and inverse-frame
congruences used in S83L3 apply with the unchanged $\sigma$.
By R82W1 and the definition of $I$, multiplication by $v^I$ puts
the target element in the image of $L_{1,c-1}$, using the tame
average as in S83L3. This avoids needing a $p^R$ lift at $c-1$,
which need not be a cut.

Form the integral source numerator by the inverse frame, at cost
$p^\sigma$. Its image in the target at the $Q$-level agrees modulo
$Q^{qM}$ with
$p^{2\sigma}Q^{h-a}(Q-p)^IP^{G+y} E_2$.
As in S83L3 divide the source numerator by $P^{G-\eta}$, since
no jump at the down level exceeds $\eta$. The coordinatewise
division is integral by distinguished division, including when
$P$ is a unit times its monic associate. The first $h-\eta$
Taylor coefficients of the new source numerator at $Q=0$ lie
in $\mathrm{Fil}^{\eta+1}=\mathrm{Fil}^a$. Its target coefficients
are those of $p^{2\sigma}Q^{h-a}(Q-p)^I P^z E_2$ there, since
$P$ is a unit at this root over the fraction field. Take the
coefficient of degree $h-a+j$ for the $j$ of F83S1. This degree
is nonnegative and below $h-\eta<qM$. It gives an integral
source direction with image $p^{2\sigma}d E_2$, $v_p(d)<t$.
Multiplying it by $p^t/d\in O$ gives exactly $p^{2\sigma+t}E_2$.
Transport by the remaining crystalline Frobenius to obtain the
middle up-filtration membership for the full arc. It is the
only up letter. The second construction of S83L1 now applies
to every original scalar, even when the constants specified at the
two down levels differ. It gives the full $\mathrm{LiftsAt}$
with integer exponent
\[
             E_1=3\sigma+t\le12U+1\le13U .
\]
Here only the generator difference of the scalar must vanish to
the prescribed order at each down level. Thus the filtered
direction suffices without making the two constants agree.
% END F83S2
